\chapter{Resumen}


\textbf{Resumen:}

En este trabajo se ha realizado el diseño conceptual y teórico, así como la construcción de un sistema semiautomático para la generación de compost a partir de desechos orgánicos a través del sistema In-Vessell que se caracteriza por tratarse de un contenedor cerrado.

En este proceso es de vital importancia tener en cuenta como es que evoluciona el ambiente al interior del contenedor, esto con el fin de mantener unas condiciones que favorezcan el crecimiento de bacterias. Estos cultivos de bacterias proveerán de diversas propiedades químicas y físicas a los residuos, transformándolos en un abono natural.

Este proyecto tiene como misión principal construir un sistema mecatrónico que materialice el trabajo realizado durante metodología de la investigación y Trabajo Terminal I. De esta manera en este texto se hará un recorrido exhaustivo de cada fase que marcó la concepción y ejecución de este proyecto. A lo largo de este trayecto se abordarán con precisión los desafíos presentados a lo largo de cada etapa de implementación.

Se examinarán las soluciones y estrategias empleadas para sortear cada uno de estos obstáculos, además se profundizará en los conocimientos adquiridos, en la evolución que el equipo presentó respecto al desarrollo del proyecto. Este análisis no solo proporcionará una visión clara de la implementación de la idea, sino que se espera también que sirva como fuente valiosa de aprendizaje para futuros proyectos.


%\textbf{Palabras Clave:} Compost, desechos orgánicos, semiautomático, In-Vessell, bacterias, monitoreo, diseño conceptual, trituración, hilos, sitio web, ambiente, aireación.\\

\textbf{Palabras Clave:} Compost, Residuos orgánicos, Sistema semiautomático, Sistema mecatrónico, Proceso In-Vessel, Crecimiento bacteriano, Contenedor cerrado.

%%%%%%fin del archivo

\endinput 